\begin{question}
This question concerns a system with three sorts of threads, with identities
of types~|A|, |B| and~|C|.  A mechanism is required to allow three threads,
one of each sort, to synchronise and to obtain one another's identities. 
To support this, we require objects that implement the following trait.
%
\begin{scala}
trait ABCSync[A,B,C]{
  def ASync(me: A): (B, C)
  def BSync(me: B): (A, C)
  def CSync(me: C): (A, B)
}
\end{scala}
%
Each |A|-thread calls |ASync|, passing in its identity; it receives back the
identities of the |B|- and |C|-threads with which it synchronises.  Likewise,
the |B|- and |C|-threads call |BSync| and |CSync|, respectively.  Each
operation will block until other threads call the other two operations.
Implement a class that extends this trait, using |Condition|s.
\end{question}

%%%%%

\begin{answer}
In problems like this, it is easiest to pick a particular order for events to
happen.  My solution follows the sequence diagram below.  Signals on |aStart|,
|bStart| and |cStart| indicate that the relevant thread can start and should
write its identity.  Signals on |bSignal| and |cSignal| indicate that the
relevant thread can read the other threads' identities. 
%
\begin{center}
% A node at #1, "write #2"
\def\nodestyle{\small\scalacolour}
\def\labelstyle{\footnotesize\scalacolour}
\def\write#1#2{
  \draw[fill] #1 ellipse (0.6pt and 2.4pt); % Note: scaling makes this a circle.
  \draw #1 node[right]{\nodestyle write \scalashape #2};
}
\def\height{10.7}
\begin{tikzpicture}[xscale = 2.8, yscale = 0.7]
\draw (0,0) node[draw] (a) {\scalastyle ASync(a)};
\draw (a) -- ++ (0,-\height);
\draw (1,0) node[draw] (b) {\scalastyle BSync(b)}; 
\draw (b) -- ++ (0,-\height);
\draw (2,0) node[draw] (c) {\scalastyle CSync(c)};
\draw (c) -- ++ (0,-\height);
\draw (3,0) node[draw] (a2) {\scalacolour next \scalastyle ASync}; 
\draw (a2) -- ++ (0,-\height);
%
\draw[<-] (a) ++ (0,-2) node[right]{\labelstyle (A1)} -- 
  node[above]{\scalashape\footnotesize aStart} ++ (-1,0); 
\write{(a) ++ (0,-2.8)}{a}
\draw[->] (a) ++ (0,-4) -- 
  node[above]{\scalashape\footnotesize bStart} ++ (1,0) 
  node[right]{\labelstyle (B1)};
\write{(b) ++ (0,-4.8)}{b};
\draw[->] (b) ++ (0,-6) --
  node[above]{\scalashape\footnotesize cStart} ++ (1,0) 
  node[right]{\labelstyle (C1)};
\write{(c) ++ (0,-6.8)}{c};
\draw[->] (c) ++ (0,-8) --
  node[above,near start]{\scalashape\footnotesize aSignal} ++ (-2,0) 
  node[left]{\labelstyle (A2)} ;
\draw[->] (a) ++ (0,-9) --
  node[above]{\scalashape\footnotesize bSignal} ++ (1,0) 
  node[right]{\labelstyle (B2)} ;
\draw[->] (b) ++ (0,-10) --
  node[above,near start]{\scalashape\footnotesize aStart} ++ (2,0);
\end{tikzpicture}
\end{center}
%
However, there are lots of other orders one could choose. 

My code is below.  The points in the code labelled (A1), (B1), (C1), (A2) and
(B2) correspond to the points with the same labels in the sequence diagram.
The corresponding signals will happen in that order, and the next
synchronisation won't start until this sequence is completed.  The
|Condition|s |aStart|, |bStart| and |cStart| need to be augmented with boolean
variables, both to block a thread until it is its turn to proceed, and, when
the thread receives a signal, to test whether it has been preempted and so has
to wait again.
%
\begin{scala}
class ABCSyncMonitor[A,B,C] extends ABCSync[A,B,C]{
  private val lock = new Lock

  /** The identities of the current (or previous) threads. */
  private var a: A = _; private var b: B = _; private var c: C = _

  /** Are the values in £a£, £b£, £c£ valid? */
  private var aFull, bFull, cFull = false

  /** Conditions on which threads wait for permission to write their values. */
  private val aStart, bStart, cStart = lock.newCondition

  /** Conditions on which threads wait to be able to read their results. */
  private val aSignal, bSignal = lock.newCondition

  def ASync(me: A) = lock.mutex{
    aStart.await(!aFull)          // (A1)
    assert(!bFull && !cFull); a = me; aFull = true
    bStart.signal()                // signal to B£-£thread at (B1)
    aSignal.await()                // (A2)
    assert(aFull && bFull && cFull)
    bSignal.signal()                 // signal to B£-£thread at (B2)
    (b,c)
  }

  def BSync(me: B) = lock.mutex{
    bStart.await(aFull && !bFull)         // (B1)
    assert(!cFull); b = me; bFull = true
    cStart.signal()                 // signal to C£-£thread at (C1)
    bSignal.await()                // (B2)
    assert(aFull && bFull && cFull)
    aFull = false; bFull = false; cFull = false
    aStart.signal()               // signal to next A£-£thread at (A1)
    (a,c)
  }

  def CSync(me: C) = lock.mutex{
    cStart.await(bFull && !cFull)         // (C1)
    assert(aFull); c = me; cFull = true
    aSignal.signal() // Signal to A£-£thread at (A2)
    (a,b)
  }
}
\end{scala}

\end{answer}

